\section{Related Work and Research Position}
\textit{Primary paper 1: URL detection and representativeness (2022).} S\'anchez-Paniagua~\textit{et al.}~\cite{sanchez2022} evaluate RF, LightGBM and other classifiers using handcrafted URL features, alongside TF-IDF/character n-grams and neural models. Their PILU-90K dataset distinguishes legitimate index and login URLs, and they examine changes across collection periods. The strongest overall reported pipeline uses TF-IDF/n-grams with logistic regression. We use their evaluation of RF/LightGBM to motivate our comparator choice, not to claim those classifiers are their proposed best method. We do not use PILU-90K, reproduce their features, or establish login-specific or temporal robustness.

\textit{Primary paper 2: cross-source generalisation (2024).} Rashid~\textit{et al.}~\cite{rashid2024} already evaluate cross-dataset transfer before proposing target-informed, unsupervised feature alignment. Therefore, unadapted transfer evaluation alone is not our contribution. They motivate our transfer question and provide the HP data source. We keep learning source-only, rebuild the HP cohort and use our own domain-disjoint splits. We neither implement UDA nor claim superior scores: the data partitions, feature representations and budgets differ. Full-text review confirms random URL splits, validation tuning, and target-distribution alignment. Their diagnostic RF feature shortlist uses impurity importance shared across dataset-trained models, whereas our ranking uses source-domain permutation evidence. Details not reported in the article are not treated as proof that a control was absent.

\textit{Adjacent feature-stability research.} Misiek and Hyla~\cite{misiek2026} study an XGBoost-based detector with cross-dataset importance criteria and computational comparisons. Their selection considers importance across datasets (Sec.~II-B), which differs from our source-domain PI ranking. Kanon~\textit{et al.}~\cite{kanon2026} likewise connect feature stability with external evaluation. Stability and cost are therefore prior work. Our source-only stage differs in the source of selection evidence and the matched-budget, bidirectional domain-disjoint comparison. DA-GAFS adopts cross-dataset importance in the spirit of~\cite{misiek2026} and adds an explicit test of directional consistency, a cost-aware budget and a cost-aware ensemble choice, each ablated. No universal priority claim is made.

\begin{table}[t]
\caption{Roles of the two primary papers and the question evaluated here. Entries distinguish scope, not comparative performance.}
\label{tab:positioning}
\centering\scriptsize
\setlength{\tabcolsep}{3pt}
\begin{tabular}{>{\raggedright\arraybackslash}p{0.22\columnwidth}>{\raggedright\arraybackslash}p{0.31\columnwidth}>{\raggedright\arraybackslash}p{0.35\columnwidth}}
\toprule
Study & Role in this project & Boundary of the comparison \\
\midrule
S.-Paniagua et al. (2022) & URL classifiers, login representativeness & PILU differs from EBBU/HP; no exact replication \\
Rashid et al. (2024) & Transfer problem, UDA, HP & Target-informed adaptation; no UDA baseline here \\
This study & Source-only ranking (frozen) and DA-GAFS (exploratory) & Matched-$k$ selectors, domain groups, fixed source thresholds, two directions; DA-GAFS uses labelled FIT of both collections \\
\bottomrule
\end{tabular}
\end{table}

\textit{Evaluation validity.} Leakage and a mismatch between the test population and the scientific claim can invalidate apparently strong results~\cite{kapoor2023}. We specify source-only nested selection, domain grouping, source-locked thresholds and bounded claims to make these assumptions inspectable. PI, RF, LightGBM and TreeExplainer are established components whose original provenance is preserved in the repository; combining them does not establish a novel learning algorithm. STABLE-PI and DA-GAFS are explicit ranking heuristics, without a transferred theoretical guarantee of feature-selection consistency.
